\section{Stopping and Depth}
\label{sec:stopping}

To shorten a hole, we first bring it into a normal form: some words are
stopped, and the others continue periodically (\cref{sec:rays}). Both parts
can then be shortened (\cref{sec:trie,sec:truncation}). This section shows
how the words of a hole can grow. Once a run is rejected, a letter either
kills it or keeps it rejected, so as long as the run stays alive the tests
at single cuts stay satisfied. The comparison tests are more sensitive: a
new letter can complete a match with a word that stays fixed.

\subsection{Joint extension}

Suppose we append a letter $c$ to $w_i$ in a hole and keep $w_j$ fixed.
The fixed target $j$ \emph{blocks} $c$ if some rest of $w_j$ equals some
rest of $w_ic$ and the pair of states at the cuts where these rests start,
the host state on $w_ic$ and the target state on $w_j$, is admitted.
Equality puts the end of the match at the new host endpoint. This includes
an empty match at the new endpoint. A word is never blocked by itself,
since the blocking word stays fixed while the blocked word grows.

\begin{lemma}[Joint extension]\label{lem:joint}
Append one letter to each coordinate in any chosen set, and keep the
other coordinates fixed. The new tuple is a hole exactly when every new
letter keeps its run alive and no fixed target blocks it.
\end{lemma}

\begin{proof}[Proof idea]
A comparison at the new end of an extended target has a live successor
of a rejected state as its target state. At every earlier cut of that
target, deleting the last letter of its rest gives an old comparison.
The host loses at most one letter and both starting states stay the same.
With a fixed target, a new comparison can only use the new host endpoint.
Those are precisely the blocked letters. \Cref{app:rays} gives the
details.
\end{proof}

\begin{figure}[t]
\centering
\begin{tikzpicture}[x=11mm,y=6.4mm,>=Stealth,
  nd/.style={draw,circle,inner sep=1.1pt,font=\small,fill=white,minimum size=6.2mm},
  rej/.style={nd,double}, letter/.style={font=\footnotesize},
  txt/.style={font=\footnotesize,anchor=east}, cut/.style={circle,fill,inner sep=1.25pt}]
  \node[txt] at (-.42,0) {host $0$};
  \node[rej] (p0) at (0,0) {$\alpha_0$};
  \node[rej] (p1) at (1.2,0) {$\alpha_1$};
  \draw[->] (p0)--node[letter,above]{$0$}(p1);
  \draw[->] (p1) edge[loop above,looseness=6] node[letter,above]{$0,1$} (p1);
  \draw[->,dashed] (p1)--node[letter,above]{$1$}(2.4,0);
  \node[cut] at (2.4,0) {};
  \node[letter,anchor=west] at (2.52,0) {proposed extension};
  \node[txt] at (-.42,-1.8) {fixed target $1$};
  \node[rej] (q0) at (1.2,-1.8) {$\beta_0$};
  \node[rej] (q1) at (2.4,-1.8) {$\beta_1$};
  \node[nd] (dd) at (3.6,-1.8) {$\dead$};
  \draw[->] (q0)--node[letter,below]{$1$}(q1);
  \draw[->] (q1)--node[letter,below]{$0,1$}(dd);
  \draw[densely dotted,thick] (p1)--(q0)
    node[letter,pos=.5,anchor=east,xshift=-1mm]{admitted pair};
\end{tikzpicture}
\caption{The reader of \cref{ex:comparable-reader} on the hole $(0,1)$.
Double circles mark rejected states; all four states of the hole are
rejected. The pair $(\alpha_1,\beta_0)$ is admitted, so appending $1$ to the
host would create an admitted comparison with the fixed target. Every letter
leads from $\beta_1$ to the dead state, so the target stops.}
\label{fig:stopping}
\end{figure}


\Cref{fig:stopping} shows the reader of \cref{ex:comparable-reader} on the
hole $(0,1)$. Keep the second word fixed. Appending $1$ to $0$ makes the
rest at cut $1$ equal to the whole word $1$, and the pair
$(\alpha_1,\beta_0)$ is admitted, so the letter $1$ is blocked. Appending
$0$ keeps the hole, and so does every further $0$. The state $\alpha_1$
hosts the blocking comparison although rejection began at cut $0$. For one
word this cannot happen (\cref{sec:comparison}). The second word stops
because either letter takes it to the dead state.

The lemma lets us choose which words to keep fixed. A coordinate
\emph{stops} relative to a set if every letter kills its run or is
blocked by a target in that set. A set of \emph{stopped coordinates}
has an order in which each stops relative to its predecessors. We call
this order its \emph{certificate}. The other coordinates are
\emph{moving}.

\begin{lemma}[Stopping]\label{lem:stopping}
Every hole has a coordinatewise extension that is a hole with a set
of stopped coordinates. From that extension all moving coordinates
can grow together by one letter at every step, with a hole at each
step and the stopped words fixed.
\end{lemma}

The proof chooses a hole extension with as many stopped coordinates as
possible. If a moving word had no live letter unblocked by the stopped
words, it could be appended to their stopping order. Joint extension
therefore permits all moving words to grow. The certificate for the
stopped words only uses those words, so it persists. In
\cref{fig:stopping} the second coordinate stops and the first one can
grow along $0^\omega$.

\subsection{Charging the depth}

The suffixes of the stopped words are the possible rests of fixed targets.
We need to know how many of them can fit along one host rest at the same
cut. This number, the depth, will bound how far along a host we must look
for a region free of stopped suffixes (\cref{sec:rays}), and it is the
measure of the induction in \cref{sec:trie}. For a \emph{test word} $x$ and a set of coordinates
$F$, define
\[
 \RestLen(W,F,x)=
 \{d:0<d\le\len x,\ \wpre xd\text{ is a suffix of some }w_i,\ i\in F\}.
\]
The largest cardinality of these sets over all test words is the
\emph{depth} of the words in $F$. This counts lengths, so multiple
occurrences of a suffix count only once.

The bound charges each length to an obstruction that can be recognized
at the start of its match. In a general reader the whole state identifies
the obstruction. In a compiled reader a single component and signed
variable will suffice.

\begin{lemma}[Depth of stopped words]\label{lem:depth}
For a hole of a reader with $N$ states and a stopped set $F$, the depth is
at most $(N+1)|F|$. For a compiled reader of a system with $v$ variables,
it is at most $(4v+1)|F|$.
\end{lemma}

\begin{figure}[t]
\centering
\begin{tikzpicture}[x=7.4mm,y=5.2mm,>=Stealth,
  lbl/.style={font=\footnotesize}, st/.style={circle,draw,fill=white,inner sep=.9pt,font=\scriptsize},
  old/.style={line width=2.2pt,black!25}, suf/.style={line width=2.2pt,black!80}]
  \draw[thick] (0,0)--(10,0);
  \node[lbl,anchor=east] at (-3.1,0) {test word $x$};
  \foreach \p/\t in {3/d,5.5/{d'},6.5/{d'+1}}
    {\draw (\p,-.18)--(\p,.18); \node[lbl,below] at (\p,-.15) {$\t$};}
  % a word whose next letter kills it
  \draw[old] (-2.2,1)--(0,1); \draw[suf] (0,1)--(3,1);
  \node[lbl,anchor=east] at (-3.1,1) {$w_i$};
  \node[st] at (0,1) {$\gamma$};
  \draw[->] (3,1)--node[lbl,above]{$x[d]$}(4.1,1); \node[lbl,anchor=west] at (4.1,1) {$\dead$};
  % a word whose next letter is blocked by an earlier word
  \draw[old] (-1.4,2)--(0,2); \draw[suf] (0,2)--(5.5,2);
  \node[lbl,anchor=east] at (-3.1,2) {$w_{i'}$};
  \node[st] at (0,2) {$\gamma'$};
  \draw[old] (-2.8,3)--(0,3); \draw[suf] (0,3)--(6.5,3);
  \node[lbl,anchor=east] at (-3.1,3) {$w_j$, earlier};
  \draw[->,dashed] (5.5,2.15)--(6.5,2.85);
  \node[lbl,anchor=west] at (6.2,2.05) {pass};
\end{tikzpicture}
\caption{Charging the lengths of $\RestLen(W,F,x)$. Each stopped word is
drawn so that its matched suffix (black) is a prefix of the test word, and
the match starts in the state at its left end. The next letter of $x$
kills $w_i$, so no longer match starts in $\gamma$. The next letter after the
match on $w_{i'}$ is blocked by a rest of an earlier word $w_j$ that starts
at or before the match, so the charge passes to the length $d'+1$ on
$w_j$.}
\label{fig:charging}
\end{figure}


To see the charge (\cref{fig:charging}), fix a test word $x$. A
\emph{match} of length $d$ is an
occurrence of $\wpre xd$ as a suffix of a word of the hole. It starts at a
cut of that word, in some state $\gamma$. Let $d<\len x$ be in
$\RestLen(W,F,x)$, with a match on a stopped word $w_i$. The next letter
$x[d]$ cannot extend $w_i$, and the certificate says why. If $x[d]$ kills
$w_i$, then reading $\wpre x{d+1}$ from $\gamma$ dies. No longer match can
start in the state $\gamma$, since it would read $\wpre x{d+1}$ from
$\gamma$ inside the hole. So $d$ is the longest match that starts in
$\gamma$. Otherwise $x[d]$ is blocked by an earlier stopped word $w_j$: a
rest of $w_j$ equals a rest of $w_i\,x[d]$, and the pair at their starts
is admitted. If this rest starts after the beginning of the match, the same
argument with admission in place of death shows that $d$ is the longest
match that starts in $\gamma$: a longer one would contain the blocking
comparison. If it starts at or before the beginning of the match, then
$\wpre x{d+1}$ is a suffix of $w_j$, and we pass the charge to the length
$d+1$ on $w_j$.

The passes go strictly earlier in the certificate. After fewer than $|F|$
passes the charge reaches $\len x$ or a length that is the longest match
starting in some state. Each length in $\RestLen(W,F,x)$ is therefore one
of these at most $N+1$ lengths minus a number of passes below $|F|$. For a
compiled reader, the state $\gamma$ can be replaced by one signed variable
in one of its two sets, a \emph{singleton}, because death and admission are
witnessed by a singleton (\cref{sec:decision}). There are $4v$ singletons.
\Cref{app:rays,app:counting} give the proofs. The same depth bound now
controls how far we must look for periodic continuations.
